\documentclass[../main.tex]{subfiles}
\begin{document}
%==========================================TIÊU ĐỀ TẬP========================================
\begin{table}[h]
    \begin{adjustwidth}{-1cm}{}
        \begin{tabular}{>{\centering\arraybackslash}p{6cm}|>{\raggedright\arraybackslash}p{12.5cm}}
            \multirow{5}{6cm}
            % Cột bên trái: ÉPISODE và số tập
            {\\[-40pt]
                \flaregothic\fontsize{20pt}{20pt}\selectfont \centering
                \color{eptype}ÉPISODE\color{black}\\
                \burbank\fontsize{72pt}{72pt}\selectfont
                \color{epnum}63\color{black}\\[6pt]
                \cafeteria\fontsize{20pt}{20pt}\selectfont\color{epnum}(5-7)\color{black}
                \\[18pt]
                \flaregothic\fontsize{14pt}{14pt}\selectfont
                \color{epattr}BIENVENUE, \colorbox{vol}{\color{Sheep} Watame} !
                \color{black}
                \\}
             &
            % Dòng thứ nhất: tên tập tiếng Nhật
            {\fotskip\fontsize{18pt}{18pt}\selectfont はじめてのおつかい【\ruby{決}{けっ}\ruby{闘}{とう}】
            }                              \\[6pt]
            % Dòng thứ hai: tên tập tiếng Anh
             & {\cafeteria\fontsize{18pt}{18pt}\selectfont The Power of Money}                                        \\[4pt]
            % Dòng thứ ba: tên tập tiếng Việt
             & {\vnmsans\fontsize{18pt}{18pt}\selectfont Thách đấu}                                                   \\[8pt]
            % Dòng thứ tư: Nhân vật xuất hiện
             & {\caviar{\fontsize{14pt}{14pt}\selectfont Nhân vật: \emp{kuon1} \emp{coco} \emp{kanata} \emp{watame}}}
            \\[8pt]
            % Dòng cuối cùng: Ngày phát hành
             & {\continuum\fontsize{16pt}{16pt}\selectfont Ngày phát hành: 19/07/2020}                                \\[18pt]
        \end{tabular}
    \end{adjustwidth}
\end{table}
%=========================================TRUYỆN CHÍNH========================================
\color{black}

% Hộp tóm tắt
\txtbx[2]{{\color{vol} \uline{Tóm tắt:}} Tsunomaki Watame - nàng cừu mới của \hll - đã có một ngày đầu tiên ``đáng nhớ'': bị kẹt trên cây, bị cả nhóm sai vặt, rồi bị một con cừu thách đấu đánh bại trong mọi cuộc thi từ vật tay, tốc độ, trí tuệ đến cả nghệ thuật. Sự xuất hiện của kẻ thách đấu bí ẩn khiến cô suy sụp và khóc trong lòng Kuon. Cuối cùng, Coco và Kanata đã trấn an cô rằng không ai có thể thay thế vị trí ``idol cừu'' của \hll. Bữa tối chào mừng kết thúc trong tiếng cười ấm áp, với một món thịt bí ẩn mà Watame vô cùng thích thú.}

\pov{Hikawa Kuon}

Tuần thứ ba. Sau cơn bão mang tên Coco và cơn giông mang tên Kanata, tôi đã nghĩ rằng mình có thể trụ vững với bất cứ điều gì. Tôi đã sai. Và tôi sẽ tiếp tục sai cho đến khi nào tôi học được bài học rằng, ở \hll, mọi thứ luôn có thể trở nên điên rồ hơn.

Vậy, người tiếp theo trong \textit{holoForce} sẽ đến văn phòng là ai nào?

Một cô cừu vàng với nụ cười hiền hòa, luôn cầm trên tay cây đàn hạc như một người bạn đồng hành không thể thiếu. Nếu Kanata là một thiên sứ đang nỗ lực để trở thành ``thiên sứ xịn'', thì cô nàng này là một ``bard hiện đại'' một nghệ sĩ lang thang trong thời đại số, mang trong mình tình yêu âm nhạc vô hạn.

Cô ấy có một tâm hồn nhạy cảm và sâu sắc, dù bề ngoài trông như một chú cừu dễ thương, nhưng thực chất lại là một nghệ sĩ thực thụ. Với giọng hát êm dịu và khả năng sáng tác tuyệt vời, cô ấy có thể khiến bất cứ ai cũng phải rung động khi lắng nghe. Tôi đã từng nghe cô ấy hát trong một buổi livestream, và tôi đã phải dừng công việc lại trong năm phút chỉ để lắng nghe. Nó khiến tôi nhớ lại những bản ballad xưa mà bố tôi hay bật.

Cô ấy luôn nở một nụ cười và mang lại bầu không khí yên bình cho mọi người, nhưng đôi khi khá trẻ con, hay bày trò nghịch ngợm với những người thân thiết. (Ơn giời không phải là với tôi, ít nhất là chưa.) Tôi nghe nói cô ấy từng giấu đàn hạc của Coco để trêu, và Coco đã mất ba tiếng mới tìm thấy nó. Tôi hy vọng cô ấy không áp dụng chiêu đó với tôi.

Cô nàng Cừu này cũng là một người khá cừ trong trò oẳn tù tì. Cũng là một người lắm mưu để thắng trò này bằng mọi giá. Tôi từng thấy cô ấy thắng liên tiếp mười ván trước Miko, và cô ấy đã cười một cách đáng sợ. Tôi không hỏi bí quyết, và tôi nghĩ tôi không muốn biết.

Và dù là một cô cừu hiền lành, mục tiêu lớn nhất của cô không chỉ là cống hiến cho âm nhạc. Cô ấy muốn một ngày nào đó có thể thực hiện được ước mơ là được đứng trên sân khấu Budokan. Một cô cừu dễ chịu và yêu thích âm nhạc.

\dots\ Hoặc có lẽ, một cô cừu lắm chiêu trò ẩn giấu dưới lớp mặt nạ hiền lành. Tôi đã học được rằng, trong công ty này, vẻ ngoài hiền lành thường là cái bẫy giỏi nhất.

\chrint

Đó là Tsunomaki Watame (\ruby{角}{つの}\ruby{巻}{まき}わため, \ruby{Sư-nô}{Giác}-\ruby{ma-ki}{Quyển}\ Oa-ta-mê), bạn thân của Coco, và cũng là thành viên thứ ba của \textit{holoForce} mà tôi sẽ phải đối mặt.

Một cô gái thích ca hát (như đã nói ở trên) và thường xuyên tải lên các bản cover (hát lại) và phát trực tiếp các buổi hát hàng tuần có tựa đề \textit{Watame Night Fever!!} – đến mức cô nàng còn có một bộ trang phục riêng cho nó, với những chiếc đèn disco nhỏ xíu đính trên áo. Trông cô ấy như một DJ cừu trong một buổi tiệc thập niên 70.

Watame có tính tình nhẹ nhàng và nói chuyện với giọng điệu dễ chịu, điềm tĩnh và lịch sự, phù hợp với hình ảnh một chú cừu mềm mại, mịn màng. Người hâm mộ nhận xét giọng nói của cô đặc biệt nhẹ nhàng, trưởng thành và nữ tính, với kiểu giọng mà nếu nghe vào buổi tối, bạn sẽ muốn uống một cốc sữa ấm và đi ngủ ngay lập tức. Cô cũng được bạn bè biết đến là người khá ấm áp và thân thiện. Nhiều khi cô ấy sẽ nói chuyện với khán giả bằng chính tiếng cừu của mình, và tôi không biết làm thế nào mà cô ấy có thể nói được hai thứ tiếng cùng lúc.

\model{watame}{24}{r}{8cm}

Trái ngược với hình ảnh hiền hòa, cô ấy cũng có một tính cách khá là\dots\ dị hợm: tự giải trí bằng cách làm những biểu cảm khuôn mặt điên cuồng. Nụ cười toe toét, loạn trí của cô ấy đã trở thành biểu tượng trên kênh của mình, thậm chí còn xuất hiện trên một số sản phẩm chính thức của chính cô ấy. Tôi đã thấy một chiếc áo phông có in nụ cười đó, và tôi vẫn chưa biết nên cảm thấy sợ hay nên mua một cái.

Ở \hll, mọi người (nhất là những thành viên săn mồi như Fubuki và Mio) đều có dã tâm ăn thịt cô ấy. Đó là điều thật sự rất lo ngại khi đôi lúc các thành viên sẽ lại trỗi dậy thú tính của mình. Tôi đã thấy Fubuki nhìn Watame với ánh mắt mà tôi chỉ thấy khi cô ấy nhìn một đĩa thịt nướng. Tôi không biết đó là sự thật hay chỉ là trò đùa, nhưng tôi không dám hỏi.

Khi chơi game, cô ấy hiếm khi tỏ ra bực bội, bởi lẽ cô ấy hay chơi những trò có tiết tấu nhẹ nhàng. Watame khá yếu về mảng kinh dị, khi mà cô ấy thường hét lên kinh hoàng trước những chuyển động nhỏ nhất trên màn hình của mình. Một lần tôi thấy cô ấy hét lên chỉ vì một cánh cửa đóng sầm trong game, và tôi đã phải tắt âm thanh để tiếp tục làm việc.

Nàng Cừu này cũng là một người cực kỳ dễ xúc động. Nếu một cảnh xúc động trong trò chơi xuất hiện, chẳng hạn như cái chết hoặc lễ cưới của một nhân vật, cô ấy sẽ bắt đầu khóc – thường phải mất một lúc để lấy lại bình tĩnh trước khi tiếp tục. Điều tréo ngoe ở chỗ, tiếng sụt sịt của cô ấy khi khóc lại giống như\dots\ tiếng xì xụp của mì ramen, vì thế cứ mỗi khi Watame chuẩn bị khóc thì nhiều người sẽ nói đùa với nhau rằng: ``đến giờ ăn mì ramen rồi.'' Tôi đã thấy điều đó xảy ra trong một buổi stream, và tôi đã không thể nhịn cười được năm phút.

Cô cũng là một người có tiếng trong lĩnh vực ASMR, như Choco. Ít nhất thì còn hơn cả trăm lần ai đó làm cho tôi mất ngủ cả đêm\dots\ Tôi vẫn nhớ lần đầu nghe ASMR của Choco, tôi đã ngủ gục trên bàn làm việc. Watame thì nhẹ nhàng hơn, nhưng cũng đủ để khiến bạn muốn cuộn tròn trong chăn và quên đi mọi lo âu.

Watame có một đôi mắt màu tím, tóc vàng mượt, dài đến lưng, thường để một bím tóc nửa buộc nửa thả biến thành một kiểu tóc đuôi ngựa ngắn và một chỏm tóc ngố ngắn, có một chiếc kẹp màu hồng ở bên trái. Trông cô ấy như một chú cừu bước ra từ một cuốn truyện cổ tích.

Bộ trang phục của Watame gồm một bộ váy ngắn màu trắng không tay, viền lông thú với một dải đen dưới ngực, một chiếc đầm ngắn màu đen xếp ly bên trong, tay áo có gân viền lông thú tách rời. Trước ngực có cài một chiếc nơ đỏ có trâm, một chiếc áo choàng viền lông thú hai mặt (với mặt trong có màu hồng), một túi thắt lưng nhỏ đeo chéo sang phải và đôi bốt viền lông thú màu đen. Cô ấy cũng mang theo một cây đàn hạc, và tôi đã thấy cô ấy dùng nó để gõ vào đầu Coco một lần, nên tôi biết rằng nó cũng có thể là một vũ khí.

\chrint

Liệu cô nàng Cừu này có làm gì để báo hại tôi hay không đây\dots\ Tôi có linh cảm rằng, với biệt danh ``con cừu'' và sự hiền lành giả tạo đó, cô ấy sẽ còn gây ra nhiều rắc rối hơn tôi tưởng. Nhưng ít nhất, cô ấy có đàn hạc, và trong một văn phòng toàn những người điên, một nghệ sĩ có thể là liều thuốc giải độc.

Hoặc là chất xúc tác cho những sự hỗn loạn. Tôi chưa thể nào biết được.

\il

Tháng Sáu ở thành phố Điện tử không hề dễ chịu. Cái nóng oi bức len lỏi vào từng ngõ ngách, khiến không khí trở nên ngột ngạt như một cái lò nướng khổng lồ. Tiếng ve kêu râm ran từ những tán cây ven đường, như một bản giao hưởng bất tận của mùa hè. Tôi ngồi trong văn phòng với cái quạt máy bật hết công suất, gió thổi tung đống giấy tờ trên bàn, nhưng vẫn chẳng đủ để xua tan cái nóng. Trong khi đó, Kanata và Coco vẫn tỏ ra bình thản. Có vẻ như một người là thiên sứ, còn người kia là Rồng thì không hề hấn gì với thời tiết kiểu này, tôi bắt đầu nghi ngờ rằng họ có điều hòa tích hợp bên trong người.

Tôi vừa mới uống một ngụm trà đá, cảm giác mát lạnh chạy dọc xuống cổ họng như một tia cứu rỗi giữa sa mạc, thì điện thoại đổ chuông. Một cái tên quen thuộc hiện lên màn hình: Watame. Tôi chưa kịp bắt máy thì đã nghe thấy tiếng hét thất thanh vọng ra từ loa: ``\textbf{Kuon-senpaaaai! Kanata!! Cứu em với!!}''

Tôi suýt phun cả trà ra ngoài. Nước trà suýt văng vào bàn phím, và tôi đã phải nhanh tay lùi lại. Bên cạnh tôi, nàng Thiên sứ đã bật chế độ cảnh giác, đôi mắt cô mở to, vầng hào quang bắt đầu sáng lên. Coco chỉ khẽ nghiêng đầu, vẻ tò mò như thể đang xem một bộ phim hài.

Tôi tặc lưỡi, bấm loa ngoài để hai người còn lại cùng nghe. Nàng Rồng từ đó cất giọng hỏi: ``Bà bình tĩnh lại nào! Chuyện gì đang xảy ra vậy? Bà đang ở đâu?''

Kanata lên tiếng, giọng đầy lo lắng: ``Watame-chan, bà có sao không? Đang ở đâu mà hét thế?''

Watame đáp, giọng run run như sắp khóc: ``\textbf{Tôi đang bị kẹt trên cây!!}''

Tôi và Kanata nhìn nhau, mắt mở to. Coco thì che miệng cười, nhưng cố tỏ ra nghiêm túc.

Tôi nhíu mày với giọng khẩn thiết ở phía đầu bên kia dây: ``Cái gì mà bị kẹt trên cây? Có cây ở đâu giữa thành phố Điện tử?''

Đầu dây bên kia không trả lời câu hỏi của tôi, mà thậm chí còn la lớn trong tuyệt vọng: ``\textbf{Giờ không phải lúc hỏi đâu! Nhanh lên đi! Em sắp rơi rồi!}''

Tôi thở dài, đặt cốc trà xuống bàn. Có vẻ như hôm nay sẽ không được nghỉ ngơi như tôi mong đợi. Tôi đứng dậy, vươn vai, và ra hiệu cho hai người kia đi theo: ``Thôi nào, đi cứu cô cừu điên rồ đó.''

\ct

Chúng tôi chạy đến địa điểm mà Watame chỉ dẫn. Dưới ánh nắng gay gắt của tháng Sáu, tôi phải nheo mắt mới nhìn rõ một cái cây khô cằn giữa lòng thành phố. Trông nó lạc lõng giữa những tòa nhà bê tông như một cột mốc của thiên nhiên đang cố gắng chống lại sự đô thị hóa. Và trên đó, một cô gái với mái tóc vàng bồng bềnh cùng đôi sừng đen đang bám chặt vào nhánh cây, mặt tái mét, hai chân run lẩy bẩy.

``Mình bị kẹt trên này rồi!''

\img{5-7a}

Không cần nhìn kỹ, tôi cũng nhận ra đó chính là Tsunomaki Watame, nàng thi sĩ cừu của \hll, người mỗi sáng vẫn thường cất giọng hát dịu dàng và đọc thơ. Nhưng hôm nay, thay vì những câu thơ trầm bổng, em ấy chỉ có thể kêu cứu, và trông em ấy như một chú cừu non bị lạc trên đỉnh núi.

Tôi khoanh tay, ngước nhìn nàng Cừu đang bám chặt lấy nhánh cây khô. Nhánh cây cong xuống dưới sức nặng của em, trông như sắp gãy bất cứ lúc nào.

``Bà leo lên đó làm gì vậy?'' Kanata hét lớn từ phía dưới, hai tay đưa lên che mắt khỏi nắng. ``Cây đó có gì đặc biệt à?''

Nàng Cừu lắp bắp, giọng vừa xấu hổ vừa sợ hãi: ``Tại tôi muốn xem nó là gì\dots\ Trông nó như một cây cổ thụ, tôi nghĩ có thể có tổ chim, hoặc quả gì đó lạ\dots''

Tôi liếc nhìn cái cây chẳng có gì đặc biệt ngoài mấy cành khẳng khiu và một vài chiếc lá héo úa: ``Có cái gì đâu mà xem? Đây là thành phố, chứ có phải đồng cỏ đâu mà tò mò? Đó chỉ là một cây khô bị ai đó quên chưa chặt thôi!''

Nàng Cừu bĩu môi, cố bào chữa: ``Lần đầu lên phố nên em mới thích khám phá chứ! Còn hơn là ngồi trong văn phòng cả ngày nhìn mấy đống giấy tờ!''

Nàng Thiên sứ nghe lời giải thích đó liền thở dài, lắc đầu, tay chống nạnh: ``Và thế là bà không cưỡng lại được sao? Một cái cây khô cũng đủ để bà leo lên?''

Watame chớp chớp mắt, vẻ vô tội, nhưng tình cảnh hiện tại chẳng giúp gì cho em ấy cả. Tôi còn chẳng biết em ấy lên kiểu gì cơ\dots\ cái cây đó trơn và cao, không có cành nào thấp để trèo.

Ba chúng tôi trao đổi ánh mắt. Cần phải nghĩ cách đưa em ấy xuống, nhưng mỗi khi chúng tôi đưa ra một phương án, nàng Cừu lại có lý do để bác bỏ ngay lập tức.

\dots Chẳng mấy chốc, cả ba chúng tôi cuốn vào một trò chơi vần điệu kỳ lạ. Tôi không biết nó bắt đầu từ lúc nào, nhưng nó đã xảy ra.

Coco khoanh tay, lên tiếng trước với vẻ mặt đầy thách thức: ``Cứ dùng móng guốc của bà là xuống được thôi!''

Watame lập tức nhăn mặt, giọng đầy phản đối: ``Dê núi (ヤギ) mới làm được thế! Tôi là cừu, không phải dê!''

Kanata thử đề xuất, tay đưa lên chỉ về phía thân cây: ``Mấy cây thân dài và hẹp thì tuột xuống dễ mà. Cứ ôm thân cây mà trượt xuống.''

Nàng Cừu lắc đầu nguầy nguậy, giọng vẫn đầy bất mãn: ``Cây hành (ネギ) mới thế! Đây là cây thật, không phải rau củ!''

Tôi nheo mắt, quan sát cái cây khô ấy một lần nữa. Vỏ cây đã bong tróc, màu nâu xám, chẳng có gì hấp dẫn. Tôi hỏi, giọng đầy tò mò: ``Mà nhắc mới nhớ, cái cây mà em đang leo là cây gì vậy?''

Watame lập tức đáp ngay, không cần suy nghĩ, như thể đã học thuộc bài: ``Nó là cây thu (はぎ)!''

Coco lập tức lấy điện thoại ra, giả vờ bấm số, miệng nói với giọng đầy kịch tính: ``A-lô? Chú cảnh sát ạ? Có một con cừu đang kêu cứu ở trên cây, nhưng nó cũng là cừu thật nên tôi không biết có nên báo cảnh sát hay không.''

Watame hoảng hốt, mặt tái hơn nữa: ``Gọi (サギ) làm cái gì vậy! Tôi không phải là kẻ lừa đảo!''

Tôi cười cười, tiếp tục trò chơi, vì tôi nhận ra đây là cách duy nhất để giữ Watame bình tĩnh: ``Bánh kem mà chúng ta ăn làm bằng gì?''

Dường như em ấy vẫn chưa nhận ra cái bẫy, liền trả lời chắc nịch như đang thi trắc nghiệm: ``Bằng bột mì (むぎ) chứ đâu!''

Kanata hứng thú tiếp tục nhập cuộc, giọng đầy vẻ của một người dẫn chương trình: ``Tên nhà ga của tuyến đường sắt Inbi phía Tây Nhật Bản là gì?''

Watame đáp ngay, không chút do dự: ``Đó là Nagi (ナギ)!''

Tôi khoanh tay, ra vẻ đăm chiêu như một nhà thám hiểm đang suy nghĩ về một câu hỏi triết học: ``Con gì có thể leo trèo dễ dàng?''

Em ấy không chần chừ, mất bình tĩnh: ``Đó là con khỉ (モンキー)!''

Coco cố nhịn cười, hỏi tiếp với vẻ thích thú: ``Thế con gì kêu `be be' nào?''

Watame dõng dạc đáp, hét thấu lên vang trời, như thể đang tuyên bố một sự thật hiển nhiên: ``Đó là con cừu (ヒツジ)!''

Em ấy bỗng khựng lại. Một giây trôi qua. Hai giây. Rồi ba giây.

Watame mở to mắt, mặt dần cứng đờ. Cô nhìn xuống chúng tôi, rồi nhìn lên cành cây, rồi lại nhìn xuống, như thể đang cố xử lý một thông tin quá tải: ``Khoan\dots\ mình là cừu mà\dots''

Chúng tôi không nhịn nổi mà bật cười. Coco phá lên, Kanata ôm bụng, còn tôi thì lắc đầu, vừa cười vừa nói: ``Nãy giờ chúng ta đang chơi vần với em ấy à? Cả một vòng tròn vần điệu mà em ấy không hề hay biết.''

Coco gật đầu, vẫn còn cười: ``Tuyệt đấy. Nhưng mà bà ấy đúng là cừu thật, tự nói ra mà không biết.''

Kanata cười khúc khích: ``Vần điệu đỉnh quá, Watame-chan tự mình nói ra cuối cùng.''

Nhận ra mình vừa bị chơi một vố đau điếng, Watame trừng mắt nhìn xuống, dang rộng hai tay ra đầy tuyệt vọng, như một nhân vật trong vở kịch bi kịch, rồi hét lên: ``\textbf{Ai thèm quan tâm chứ!? Cứu tôi đi!!!}''

Tôi nhìn em ấy, rồi nhìn Kanata và Coco, thở dài: ``Rồi, rồi, không đùa nữa\dots\ Phải thực sự đưa cô ấy xuống dưới thôi.''

Ba chúng tôi chụm đầu lại, nghiêm túc suy nghĩ. Tôi nhìn lên nàng Cừu vẫn đang bám chặt trên nhánh cây, trông em ấy như một con thú bị dồn vào chân tường, mắt mở to, tay chân run run - rồi quay sang hai người còn lại: ``Hai người có cao kiến gì không?''

Nàng Thiên sứ khoanh tay, lắc đầu, vầng hào quang trên đầu khẽ rung lên như đang suy nghĩ: ``Cái đó thì em chịu\dots\ Bà thì sao, Coco?''

Nàng Rồng cười tự tin, vỗ ngực một cái bốp, giọng đầy khí thế: ``Cứ thế này mà triển thôi! Tôi có một kế hoạch nhanh gọn, hiệu quả, và đậm chất Rồng.''

Em ấy bày ra một kế hoạch mà tôi gọi là ``triệt để'': lấy một cái rìu chặt phăng cái gốc cây xuống. Với sức mạnh của cô ấy, việc này không phải là vấn đề. Trong tưởng tượng của em ấy, Watame thậm chí còn reo lên vui vẻ, hai tay giơ lên trời như một nữ hoàng được giải cứu.

\img{5-7c}

Nhưng ngay lập tức, Watame ngoài đời thực tái xanh mặt, thét lên hoảng loạn, giọng như tiếng còi báo động: ``\textbf{Tôi chưa muốn chết đâu! Tôi còn chưa hát xong bài hát của mình!}''

Tôi thở dài. Cách này có vẻ hơi\dots\ triệt để quá. Đành nghĩ ra một phương án khác vậy: ``Hoặc là ba bọn anh làm tấm nệm đỡ em xuống. Em cứ việc nhảy, bọn anh sẽ đỡ cho! Chúng tôi đã từng đỡ được Coco khi cô ấy nhảy từ trên kệ xuống, nên em cũng sẽ ổn thôi.''

Nàng Cừu lập tức lắc đầu nguầy nguậy, ánh mắt đầy nghi ngờ, như thể đang nhìn một lũ người phát điên: ``\textbf{Không, không, không! Mấy anh không đáng tin tí nào! Coco thì có thể đỡ, nhưng Kanata thì có ngã xuống người em còn đau hơn ngã từ trên cây!}''

Tôi nhún vai, đưa tay che miệng rồi thì thầm vào tai Coco, giọng đầy âm mưu: ``\hl{Chặt cụ cái cây luôn đi. Không thích nhẹ nhàng thì ta chơi bài nặng tay vậy. Dù sao cây cũng khô, chặt cái là đổ.}''

Coco ngay lập tức gật đầu, đôi mắt sáng lên như thể vừa nhận được một nhiệm vụ quan trọng. Em ấy chạy đi tìm rìu, bước chân nhanh nhẹn, trông như một chiến binh Viking chuẩn bị ra trận. Tôi liếc sang Kanata, thấy cô ấy cũng đã chuẩn bị tinh thần, tay nắm chặt, chỉ chờ thời cơ ra tay.

``N-Này! Mấy người đang làm gì vậy!?'' Watame hốt hoảng, giọng như sắp khóc, nhưng chúng tôi không trả lời.

Coco cầm được rìu rồi, quay lại với một nụ cười đầy tự tin. Tôi hít sâu một hơi, giơ tay ra hiệu: ``Hai, ba\dots''

Cả ba chúng tôi đồng thanh hô, giọng vang vọng giữa phố xá: ``\textbf{Chặt!}''

``\textbf{Khoan, từ từ\dots!}''

Lưỡi rìu vung xuống. Tiếng gỗ rạn nứt vang lên, và ngay lập tức, cái cây bắt đầu nghiêng dần. Những nhánh cây khô gãy răng rắc, lá khô rơi xuống như mưa.

``\textbf{Waaaaaaa!!!}'' nàng Cừu hét lên thất thanh, tay vẫn bám chặt vào thân cây, mắt nhắm nghiền.

Khi cái cây bắt đầu đổ xuống, tôi thấy nàng Cừu mở to mắt, rồi bỗng nhiên trông như\dots\ ngộ ra điều gì đó. Giây phút đối diện với ``cái chết'' (trong suy nghĩ của em ấy) khiến Watame chiêm nghiệm lại cả cuộc đời mình.

\img{5-7d}

\begin{center}
    \uline{Hồi tưởng\\15 phút trước.}
\end{center}

Ngay khi trở thành ``lính mới'' trong văn phòng, cô ấy đã bị đàn chị cùng với đồng nghiệp cùng thế hệ sai vặt không thương tiếc. Một loạt yêu cầu được đưa ra như một bản danh sách mua sắm dài ngoằng:

Matsuri  khoanh tay, cười ranh mãnh: ``Watame, mua hộ chị ít đồ ăn vặt đi! Chị thèm snack vị mực nướng.''

Suisei chống cằm, nói với vẻ bình thản như đang đọc danh sách việc cần làm: ``Mua nước ngọt cho chị nhé, Watame-chan. Không đường, có gas, lạnh.''

Flare thì chậm rãi nhắc nhở, giọng đầy vẻ ``chị cả'': ``Mua cho chị cái gì để ăn tối nhé, Watame-chan! Đồ gì nhẹ nhẹ thôi, đỡ phải tốn công nấu.''

Tất nhiên, tôi cũng không ngoại lệ, tôi vỗ vai em ấy, giọng đầy triết lý: ``Em mua hộ anh kẹo cao su ở cửa hàng tiện lợi ha, Cừu-chan. Lấy cái hương bạc hà, không đường.''

Cuối cùng, Kanata nheo mắt, vỗ vai nàng Cừu với một vẻ mặt bí ẩn: ``Đi mua gì đó đi nha.''

Tôi nhìn sang, thấy Watame đứng đơ như tượng, hai tay cầm một mảnh giấy ghi chép nhưng không có gì trên đó. Tôi nhún vai, toát mồ hôi: ``Mơ hồ thế em? Anh thấy mặt em như vừa bị yêu cầu giải một bài toán tích phân ấy.''

\begin{center}
    \uline{Kết thúc hồi tưởng.}
\end{center}

Tiếng rít của gió kéo tôi về thực tại. Ngay khi cái cây bắt đầu đổ xuống, nàng Cừu đột nhiên bật nhảy lên cao, xoay một vòng trên không, tay duỗi thẳng, chân co gọn, rồi tiếp đất bằng một tư thế hoàn hảo: một chân quỳ, một tay chống đất, trông như một vận động viên thể dục dụng cụ vừa hoàn thành bài thi Olympic.

\gif{5-7e}{1}{105}

Cả ba chúng tôi đứng hình. Coco đứng yên với cây rìu trên tay, Kanata há hốc miệng, còn tôi thì quên cả thở.

Watame thở hổn hển, hai tay chống gối. Sau một giây định thần, em ấy ngẩng lên, trừng mắt nhìn Kanata, giọng đầy phẫn nộ, như một người vừa trải qua một cơn ác mộng: ``Ít nhất thì\dots\ cũng phải nói\dots\ là bà muốn gì chứ, đồ dở hơi kia!!! Tôi đã phải leo cây, bị chặt cây, suýt chết, chỉ vì bà không nói rõ là muốn ăn gì!''

Chúng tôi đưa mắt nhìn nhau, không thể tin nổi vào những gì vừa thấy. Watame không lẽ vừa thực hiện một loạt các chuỗi hành động cực khó và rồi tiếp đất hoàn hảo như vận động viên thể dục dụng cụ!? Đôi mắt tôi mở to, và tôi nhìn em ấy như thể em ấy vừa mọc thêm một đôi cánh khác.

\dots Hoặc có lẽ do em ấy uất ức gì đó với Kanata nên em ấy chưa chết ngay được. Tôi đoán vậy. Uất ức có thể là động lực mạnh hơn cả adrenaline.

Tôi lặng lẽ áp sát vào tai Kanata, thì thầm với giọng đầy sự nhẹ nhõm: ``\hl{Lần sau nói rõ cho người ta biết mình muốn gì đi chứ. Nói như vừa rồi người ta cáu là phải!}

Kanata không đáp, chỉ gật gù ra vẻ đã hiểu. Sau đó, em ấy quay sang nhìn Watame vẫn còn đang tức giận, bình thản nói như chưa có gì xảy ra: ``FamiChicken nhé! Có vẻ bà đã sẵn sàng để đi mua rồi.''

Watame chỉ biết há hốc miệng, rồi cúi đầu thở dài, như một người đã đầu hàng số phận.

\ct

Watame bước ra khỏi cửa hàng tiện lợi với khuôn mặt cau có, tay xách lỉnh kỉnh đủ thứ đồ mà chúng tôi đã sai em ấy đi mua. Trông em ấy như một cái cây giáng sinh di động, với túi này treo trên tay, túi kia vắt vai, và một chiếc túi nylon màu trắng kẹp dưới nách.

\img{5-7f}

Ngẫm nghĩ lại thì, quả thật tội cho em ấy thật. Mới bước chân vào văn phòng đã bị cả đám sai vặt như một đứa em út thời đại học. Cũng tại em ấy hiền quá mà. Nếu là Coco, chắc em ấy đã bảo ``tự đi mà mua, đừng có sai tôi'' từ ván đầu.

Tôi và hai người còn lại đứng đợi ở ngoài cửa, mặc cho cái nắng tháng Bảy chiếu thẳng xuống đầu. Kanata reo lên phấn khích khi thấy gói FamiChicken trong tay Watame, tay cô giơ lên như thể vừa nhìn thấy báu vật quốc gia: ``Tuyệt! FamiChicken! Bà đúng là thiên thần! Tôi yêu bà!''

Nhưng chưa kịp để Kanata vui mừng trọn vẹn, bỗng Watame đột ngột dừng bước. Tôi thấy em ấy khẽ nhíu mày, ánh mắt trở nên cảnh giác, hai tai cừu khẽ dựng lên như một radar động vật đang quét sóng.

``\textbf{Ai đó!?}''

Em ấy bất ngờ hét lên, xoay người sang bên phải. Tôi lập tức nhìn theo, không hiểu chuyện gì xảy ra.

Trước mặt chúng tôi là\dots\ một con cừu. Một con cừu bình thường, nhưng trông nó có gì đó khác lạ. Nó đứng thẳng, hai chân sau vững vàng trên mặt đất, và trước ngực nó đang kéo lê một phong thư màu nâu với dòng chữ nguệch ngoạc: ``Thư thách đấu.''

\img{5-7g}

Tôi chớp mắt, nhìn lại lần nữa để chắc chắn mình không nhìn nhầm. Một con cừu đang giao thư? Giữa lòng thành phố Điện tử? Chuyện này còn kỳ lạ hơn cả việc Watame leo cây.

Coco hít sâu một hơi, ánh mắt như thể vừa tìm ra một kho báu bí ẩn. Cô nghiêng đầu, nhìn con cừu với vẻ thích thú: ``Không lẽ\dots\ đây là thư thách đấu mà tôi đã nghe đồn?''

Tôi nhớ ra thứ này - bức thư mà tôi đã thấy từ tuần trước, với ba chữ ``Thư thách đấu'' và hình con cừu. Hóa ra nó không phải là trò đùa của ai đó. Nó là sự thật.

Cả hai cùng thốt lên: ``Không lẽ nào\dots?''

Kanata tròn mắt, chỉ tay vào con cừu thường kia, giọng đầy phấn khích: ``Nó đang muốn khiêu chiến với Watame để giành ghế làm idol cừu của \hll!? Một cuộc chiến sinh tồn giữa các loài cừu!''

Ngay lập tức, con cừu đó\dots\ gật đầu. Nó gật một cái rõ ràng, dứt khoát, như thể nó hiểu chúng tôi đang nói gì.

Tôi đơ người mất một giây. Một con cừu gật đầu. Nó hiểu tiếng người. Coco cũng trố mắt ra nhìn, miệng há hốc. Kanata thì\dots\ chắc cũng như chúng tôi, vừa sốc vừa tò mò.

Nhưng trong khi chúng tôi còn đang sốc, thì người đáng lẽ phải ngạc nhiên nhất lại lắc đầu nguầy nguậy, chéo hai cánh tay tạo hình chữ X, giống như một tín hiệu giao thông báo ``cấm'': ``Khỏi đi. Tôi không có hứng.''

Lời từ chối dứt khoát của em ấy như sét đánh ngang tai của kẻ thách đấu. Con cừu đối diện đứng đơ ra, mắt tròn xoe, trông như thể vừa bị ai đó tát vào mặt. Nó đứng bất động vài giây, rồi hai chân bắt đầu run lên.

\begin{figure}[H]
    \centering
    \begin{subfigure}[t]{0.45\textwidth}
        \centering
        \includegraphics[width=8cm]{../../../../Image/Book5/5-7h1.png}
        \caption{Watame đã nói ``không''\dots}
    \end{subfigure}
    \quad
    \begin{subfigure}[t]{0.45\textwidth}
        \centering
        \includegraphics[width=8cm]{../../../../Image/Book5/5-7h2.png}
        \caption{\dots\ làm cho con cừu - kẻ thách đấu - bị sốc nặng.}
    \end{subfigure}
\end{figure}

Đã thế, tôi còn buột miệng nói thêm một câu, với giọng đầy vẻ khinh khỉnh: ``Thật. Một con cừu bình thường thế này thì làm được trò trống gì? Nó còn chẳng biết nói tiếng người, trừ khi nó biết gật đầu thì thôi.''

Như một nhát dao chí mạng, con cừu kia cúi đầu ủ rũ, hai tai cụp xuống, trông như thể vừa mất hết ý chí sống. Nó nhìn xuống đất, như một kẻ vừa bị từ chối tình cảm trong một bộ phim lãng mạn.

Kanata hoảng hốt, ngay lập tức chạy lại vỗ về con cừu, giọng đầy thương cảm: ``Tại sao vậy!? Nó đã cất công đến tận thành phố của con người này kia mà! Trông nó như vừa đi bộ từ trang trại đến đây ấy!''

Coco thì đập mạnh vào lưng Watame, giọng đầy khích lệ: ``Mau chấp nhận lời khiêu chiến của nó đi! Đừng làm một con cừu đồng loại phải thất vọng chứ! Đó là danh dự của loài Cừu các bà mà!''

\img{5-7i}

Tuy vậy, Watame vẫn khoanh tay, bĩu môi, giọng đầy kiêu ngạo: ``Tôi không muốn phí thời gian với một con cừu hạ đẳng như thế đâu. Nó có biết hát không? Có biết chơi đàn hạc không? Tôi nghi ngờ nó thậm chí còn không biết nhai cỏ đúng cách.''

Tôi cũng gật đầu đồng tình, nhìn con cừu với vẻ đầy thương hại: ``Với cả anh nghĩ con cừu phèn này sẽ có tài cán gì đâu. Bỏ đi. Chắc nó chỉ biết ăn cỏ và ủ rũ thôi.''

Ngay sau đó, con cừu kia run rẩy\dots\ rồi lôi từ một túi nhỏ trên lưng ra một tờ mười nghìn yên. Tờ tiền màu xanh lá cây sáng bóng trong ánh nắng.

*SOẠT!*

\img{5-7j}

Watame lập tức lật mặt nhanh như chớp, mắt sáng rực lên như hai ngọn đèn pha. Em ấy nắm lấy phong thư, tay giật ngay tờ tiền từ con cừu, rồi giơ nắm đấm lên cao, hào hứng tuyên bố, giọng như một chiến binh sẵn sàng ra trận: ``Ngay và luôn nào! Tôi chấp nhận thách đấu! Mười nghìn yên là đủ để tôi bỏ qua mọi tủi nhục!''

Tôi đứng hình. Coco nhìn tôi, tôi nhìn lại Coco. Kanata thì há hốc miệng, hai tay bất động giữa không trung, không biết nói thế nào.

Tôi chớp mắt vài lần, rồi lẩm bẩm: ``\vie{Vãi chưởng\dots} Chỉ mười nghìn yên đã đủ để lật kèo. Em ấy đáng lẽ nên đi làm diễn viên hài, chứ đừng làm idol cừu.''

Con cừu kia, sau khi bị giật tiền, chỉ biết đứng đó với vẻ mặt hài lòng, dù trông nó vẫn như một con cừu bình thường, nhưng tôi có thể thấy một nụ cười nở trên khuôn mặt lông xù của nó.

\ct

Tôi chưa bao giờ nghĩ rằng sẽ có một ngày phải chứng kiến cảnh tượng này.

Tsunomaki Watame - thần tượng cừu của \hll, kẻ có danh hiệu ``Wata no Uta'' vang danh bốn bể, người mà mỗi tối thường cất giọng hát làm cả công ty phải lặng im lắng nghe - hiện tại lại đang bị một con cừu thường xuyên nhuyễn hết trận này đến trận khác. Một con cừu bình thường, không tên tuổi, có vẻ như vừa mới từ một trang trại nào đó đi bộ đến thành phố.

\begin{figure}[H]
    \centering
    \begin{subfigure}[t]{0.45\textwidth}
        \centering
        \includegraphics[width=8cm]{../../../../Image/Book5/5-7k1.png}
    \end{subfigure}
    \quad
    \begin{subfigure}[t]{0.45\textwidth}
        \centering
        \includegraphics[width=8cm]{../../../../Image/Book5/5-7k2.png}
    \end{subfigure}
    \quad
    \begin{subfigure}[t]{0.45\textwidth}
        \centering
        \includegraphics[width=8cm]{../../../../Image/Book5/5-7k3.png}
    \end{subfigure}
    \quad
    \begin{subfigure}[t]{0.45\textwidth}
        \centering
        \includegraphics[width=8cm]{../../../../Image/Book5/5-7k4.png}
    \end{subfigure}
\end{figure}

Mọi chuyện bắt đầu với phần thi sức mạnh.

Tôi, Coco và Kanata đứng xem với sự tò mò, trong khi hai kẻ thách đấu đặt tay lên bàn, đối diện với nhau. Watame đã lau sạch mồ hôi, hít một hơi thật sâu, sẵn sàng cho trận chiến. Tôi còn nghĩ rằng dù sao đi nữa, Watame cũng là người - có hai tay, có sức nặng, có trọng lực - chắc chắn khỏe hơn một con cừu bình thường chứ?

Đúng không?

``\textbf{Bắt đầu!}'' Coco ra hiệu, giọng đầy hào hứng.

Chưa đầy một giây sau\dots

*RẦM!*

Tay của nàng Cừu bị ghì xuống mặt bàn nhanh đến mức không ai kịp phản ứng. Mặt bàn rung lên như có một cú va chạm nhẹ, và Watame chỉ kịp nhìn xuống bàn tay mình đã bị đè bẹp.

Tôi chớp mắt. Cả hai cô nàng đồng nghiệp với em ấy cũng chớp mắt.

``\dots\ Ể?''

Watame đờ đẫn nhìn xuống bàn tay mình, như thể vừa bị một cú đấm từ một lực lượng siêu nhiên. Trong khi đó, con cừu kia chỉ điềm nhiên lắc lắc đầu, như thể vừa làm một chuyện quá đỗi bình thường, như thể nó vừa ăn một bữa cỏ và đánh bại một idol chỉ trong vài giây.

Tôi khẽ ho một tiếng, cố gắng giữ vẻ bình tĩnh: ``À\dots\ một trận thôi chưa nói lên điều gì. Tiếp tục nào.''

Nhưng rồi mọi chuyện trở nên tệ hơn.

Phần thi tài năng, Watame bị đánh bại bởi một con cừu biết trình diễn âm nhạc, ngay trên sân nhà của em ấy. Con cừu thách đấu lấy ra một chiếc kèn harmonica nhỏ và thổi một giai điệu hoàn hảo, trong khi Watame, với cây đàn hạc của mình, chỉ kịp đánh một nốt sai trước khi bị phán ``người thắng là con cừu''. Em ấy nhìn cây đàn hạc như thể nó vừa phản bội em.

Phần thi tốc độ, con cừu đó vượt qua Watame trong một cuộc đua ngắn mà chênh lệch còn xa đến mức tôi tưởng nó đang bật chế độ gian lận chạy nhanh. Nó phóng như một mũi tên, bỏ xa cô nàng ngay từ vạch xuất phát. Watame chỉ kịp thở hổn hển, tay chống gối, nhìn theo cái đuôi lông xù biến mất ở vạch đích.

Phần thi trí tuệ, Watame cố gắng giải một bài toán đơn giản: ``2 + 2 = ?'', và con cừu kia dùng chân gõ lên bảng con, đưa ra đáp án ``4'' trước cả khi em ấy viết xong câu hỏi. Cô nàng nhìn bảng, rồi nhìn con cừu, như thể vừa bị một con vật chứng minh rằng mình không thông minh bằng.

Cả ba chúng tôi đều đơ người ra nhìn cảnh tượng trước mắt. Một nàng Cừu đang bị một con cừu thường không hơn không kém vượt mặt trên mọi phương diện.

``Không thể trông mặt mà bắt hình dong được,'' tôi lẩm bẩm, đưa tay xoa trán.

Nàng Cừu thì tái mét mặt, mồ hôi chảy ròng ròng. Tôi có thể thấy rõ cảm giác sợ hãi đang bủa vây lấy em ấy, như thể em vừa nhận ra mình đang đứng trước một kẻ thù không thể đánh bại.

Watame lẩm bẩm trong cơn hoảng loạn, giọng như một linh hồn đã từ bỏ hy vọng: ``P-Phải làm sao đây\dots\ Con cừu đó không phải loại tầm thường! Nó có thể là một chiến binh cổ đại cải trang, hoặc một sinh vật từ chiều không gian khác!''

Tôi thì cũng bắt đầu thấy lo lo. Quái lạ thật, đây rõ ràng chỉ là một con cừu bình thường, nhưng nó lại khiến Watame thua như chơi. Tôi nhìn con cừu một lượt, từ đầu đến chân, không thấy gì đặc biệt ngoài bộ lông trắng bông và cặp sừng cuộn xinh xắn.

``Trông nó ghê vãi lềnh\dots''  tôi nghĩ thầm, theo bản năng lùi một bước.

Watame bắt đầu run rẩy. Có lẽ em ấy cũng đã mường tượng ra một viễn cảnh không mấy tươi sáng, nơi em bị một con cừu thường giành mất ngôi vị idol cừu: ``Nếu cứ thế này thì\dots\ nó sẽ cướp mất ngôi của mình mất! Rồi mình sẽ phải về trang trại sống!''

Trong lúc em ấy còn đang lo sợ với số phận của mình, bỗng từ phía sau có hai giọng nói vang lên: ``Có lẽ bà cứ nên nhận thua đi.''

``Hả!?'' Watame giật bắn người, quay lại nhìn chúng tôi.

Đó là tiếng của Coco, người mới cất tiếng khuyên bảo với một vẻ mặt đầy thông cảm: ``Thua cũng không có gì xấu. Dù sao bà vẫn là Watame mà, không cần phải chứng minh gì cả.''

Kanata cũng lên tiếng, giọng dịu dàng như một người chị đang an ủi em út: ``Đúng đó, bà đã có bọn này rồi, phải không nào? Không cần phải cạnh tranh với một con cừu lạ kia nữa.''

Tôi cũng khẽ mỉm cười, đặt tay lên vai em ấy: ``Còn có cả anh nữa mà! Mười nghìn yên hay không, em vẫn là Watame-chan. Và bọn anh vẫn ở đây mà.''

Watame chớp mắt, chưa kịp phản ứng gì thì nàng Rồng tiếp lời, giọng đầy tự hào: ``Bọn tôi biết rõ ai mới là idol cừu của \hll\ mà. Không phải con cừu thường kia đâu.''

Nàng Thiên sứ gật đầu, vỗ vỗ vào đầu con cừu thách đấu: ``Phải không nào, Watame?''

\img{5-7l}

Nhưng giây tiếp theo, tôi nghe thấy một tiếng kêu đầy uất ức vang lên từ phía Watame: ``\textbf{Lộn hàng rồi má / mấy đứa!!}''  cả tôi và Watame đồng thanh thốt lên khi thấy cảnh tượng trước mắt.

Thay vì đi an ủi Watame của chúng ta, Coco và Kanata lại đang vui vẻ với con cừu thường này. Kanata đang vuốt ve bộ lông mềm mại của nó, còn Coco thì mỉm cười nhìn nó đầy trìu mến, như thể vừa tìm được một người bạn mới. Hai đứa này rốt cuộc có biết mình đang cổ vũ cho ai không vậy!?

Tôi nhìn sang Watame. Em ấy mở to mắt, nhìn hai người bạn của mình bằng ánh mắt tổn thương sâu sắc, như thể vừa bị đồng đội phản bội ngay trong trận chiến cuối cùng.

Rồi bất thình lình\dots

``\textbf{Waaaaaaaah!!}''

Nàng Cừu òa khóc, nhào thẳng vào lòng tôi. Cú va chạm bất ngờ khiến tôi lùi một bước, suýt ngã, nhưng rồi tôi nhanh chóng lấy lại thăng bằng, ôm lấy em ấy, nhẹ nhàng đặt tay lên đầu em, vỗ về như dỗ dành một đứa trẻ vừa mất đồ chơi yêu thích.

``\dots\ Ít nhất thì em vẫn còn có anh ở đây,'' tôi thở dài, nhìn về phía con cừu đối diện. ``Chả biết chúng ta sẽ làm gì được nữa\dots\ nhưng đừng lo, anh không bỏ em đâu.''

Watame không nói gì, chỉ vùi mặt vào ngực tôi, run rẩy và tuôn nước mắt như suối. Tôi có thể cảm nhận được hơi thở nặng nề của em, sự lo lắng và buồn bã vẫn chưa tan biến. Đôi tay em nắm chặt áo tôi, như thể sợ rằng tôi cũng sẽ bỏ em mà đi.

Rồi từ từ, nhịp thở của em ấy trở nên đều đặn hơn. Tôi nhìn khuôn mặt nhem nhuốc lệ ấy, và rồi khẽ ngạc nhiên.

Em ấy\dots\ đã ngủ thiếp đi từ khi nào. Đôi mắt nhắm nghiền, khuôn mặt bình thản, không còn vẻ sợ hãi hay lo lắng. Trông em như một chú cừu nhỏ vừa tìm thấy một bến đỗ an toàn.

Tôi nhìn con cừu thách đấu, rồi nhìn Coco và Kanata đang chơi đùa với nó, rồi nhìn Watame đang ngủ trong tay tôi, và thở dài: ``\textit{Ít nhất thì có một con cừu đã tìm thấy bình yên. Dù cho con cừu còn lại vẫn đang bị phản bội.}``

%==========================================TRUYỆN PHỤ=========================================
\omk

Sau khoảng ba tiếng đồng hồ chìm trong giấc ngủ, Watame từ từ mở mắt. Mơ màng, cô cảm nhận được sự ấm áp và mềm mại bao quanh: một chiếc chăn dày được đắp lên người. Khi nhận thức rõ ràng hơn, cô mới thấy Coco và Kanata đang ngồi bên cạnh, ánh mắt của họ đầy vẻ vui vẻ và nhẹ nhõm, như thể vừa trải qua một điều gì đó thú vị.

Điều kỳ lạ hơn nữa, con cừu thách đấu đã không còn ở đây. Không một dấu vết nào của nó. Như thể nó chưa hề tồn tại.

Coco nghiêng đầu nhìn cô nàng vừa ngủ dậy, nở một nụ cười đầy trêu chọc: ``Ngủ ngon chứ, Đùi cừu?''

Kanata cũng chêm vào, giọng đầy vẻ thích thú: ``Thấy bà ngủ trông dễ thương ghê luôn ấy. Ngáy cũng dễ thương, mà còn kêu be be.''

Watame ngơ ngác nhìn quanh phòng một lượt. Mọi thứ vẫn y nguyên, không có dấu vết gì của cuộc thi đấu khi nãy. Không có con cừu, không có bảng điểm, không có tờ tiền mười nghìn yên nào cả.

Cô chớp mắt vài lần, rồi thắc mắc: ``Tôi\dots\ vừa nằm mơ à? Hay là tôi vừa bị ảo giác? Mùa hè nóng nực quá nên\dots''

Coco nghiêng đầu, trông có vẻ thích thú: ``Không biết được. Bà thử kể cho bọn tôi giấc mơ của bà nào? Bà đã mơ thấy gì?''

Watame vội vàng thuật lại toàn bộ diễn biến cuộc thi kỳ lạ mà cô đã trải qua: từ trận vật tay thảm hại, cuộc đua bị thua xa, biểu diễn nghệ thuật bị đánh bại trên sân nhà, cho đến cả màn so tài trí tuệ với con cừu kia. Cô kể với giọng đầy phẫn uất, tay vẽ vẽ lên không trung như một nhân vật trong một vở kịch.

Vừa dứt lời, cô liền thấy hai người bạn của mình bụm miệng cười, vai run lên bần bật. Coco đập tay xuống đùi, còn Kanata thì gục đầu xuống thành ghế.

``Hai người cười cái gì chứ!?'' Watame bực tức, giậm chân xuống sàn. ``Tôi sợ đến chết thế này mà hai bà còn cười!? Tôi đã tưởng mình sẽ mất vị trí idol cừu đấy!''

Nàng Rồng vừa ôm bụng cười vừa phẩy tay, lau nước mắt: ``Xin lỗi, xin lỗi! Nhưng mà câu chuyện của bà nghe buồn cười lắm luôn. Một con cừu thách đấu, lại còn biết vật tay, biết đánh đàn, biết chạy nhanh\dots\ nghe như phim hoạt hình ấy.''

Nàng Thiên sứ gật đầu, khuôn mặt vẫn chưa hết nét thích thú: ``Cả công ty \hll\ này chỉ có một con cừu thôi. Không có con nào khác, trừ khi có con cừu lạc vào văn phòng từ một trang trại nào đó.''

Watame chưa kịp hiểu ý hai người bạn thì tôi đã đứng tựa vào cửa bếp, cười hiền hòa nhìn cô. Tôi bước tới, nhẹ nhàng xoa đầu cô, giọng đầy trìu mến: ``Sẽ chẳng có ai thay thế được vị trí idol cừu của em cả, Watame-chan. Cho dù con cừu đó có giỏi đến đâu, nó cũng không phải là em.''

Kanata cũng vỗ nhẹ lên vai cô, mỉm cười rạng rỡ: ``Chỉ có bà mới có thể đứng ở vị trí đó thôi! Bà là Watame của chúng tôi. Cả công ty đều biết điều đó.''

Coco khoanh tay, gật đầu đầy tự hào: ``Đúng đó, đúng đó\~{} Dù bà có thua một con cừu nào đó, thì bà vẫn là idol của bọn tôi. Bọn tôi không cần một idol cừu khác.''

Watame nghe cả ba chúng tôi khẳng định chắc nịch như vậy mà lòng chợt dâng lên một cảm giác ấm áp lạ thường. Cô mở to mắt, nhìn từng người một - Coco với nụ cười tự tin, Kanata với ánh mắt dịu dàng, và tôi với bàn tay đang đặt trên đầu cô.

Đến lúc này, cô mới thực sự nhận ra rằng mình chưa từng đơn độc. Những giọt nước mắt lại dâng lên trong khóe mắt cô, nhưng lần này, cô nhanh chóng lau đi và thay bằng một nụ cười nhẹ nhõm, thoải mái. Cô lao vào ôm cả bọn chúng tôi: ``Mọi người à\dots!''

Không ai nói gì thêm. Họ chỉ lặng lẽ trao cho nhau những cái ôm thật chặt - những cái ôm kéo dài cả phút, chẳng ai muốn rời ai. Như thể đây không chỉ đơn thuần là một hành động an ủi, mà còn là một minh chứng cho tình bạn bền chặt giữa họ. Tôi cũng siết nhẹ vòng tay, mỉm cười hài lòng khi cảm nhận được sự tin tưởng và gắn kết của nhóm.

Trong lòng tôi khẽ vang lên một suy nghĩ: ``\textit{Tao rất tiếc, cừu non ạ. Mày rất tài năng, rất giỏi, rất nhanh nhẹn, rất thông minh\dots\ Nhưng cuối cùng thì, idol cừu của chúng ta chỉ có một mà thôi - và em ấy là người duy nhất. Không thể nào khác được cả.}'' Tôi khẽ thở dài, nhưng cũng không thể giấu nổi nụ cười.

Trong bếp, lò nướng vẫn đang đỏ lửa, tỏa ra một mùi thơm hấp dẫn. Món thịt mà tôi đã chuẩn bị từ trước đã sẵn sàng cho bữa tối hôm nay - một bữa tiệc chào mừng thành viên mới của đại gia đình \hll\ tuy náo loạn mà đầy ắp tình cảm này.

\ct

Bữa tối ngày hôm đó, tôi và ba thành viên của Thế hệ thứ Tư đã có một bữa tiệc nho nhỏ. Không khí ấm áp, tiếng cười rộn rã, và mùi thịt nướng lan tỏa khắp phòng.

Watame cầm lấy một miếng đùi, cắn một miếng lớn, rồi rên lên trong sung sướng, giọng đầy thỏa mãn: ``Món này ngon thật á\dots\ Mà, Kuon-senpai, anh lấy loại thịt này ở đâu nhỉ? Em chưa từng ăn thịt nào ngon thế này!''

Tôi khẽ nhếch mép, nhưng không nói gì. Trong đầu tôi thoáng hiện một suy nghĩ mà tôi không dám thốt ra: ``\textit{Anh ngạc nhiên là em còn có thể ăn thịt đồng loại được đấy\dots\ Nhưng thôi, kệ đi, miễn là em vui.}''

Coco nhoẻn miệng cười, tay cầm một cục thịt khác, giọng đầy ẩn ý: ``Thì, ai mà biết được dâu, nhỉ? Chỉ cần nó ngon là được. Đừng hỏi nhiều, kẻo mất ngon.''

Kanata cười mỉm, gắp thêm một miếng thịt bỏ vào bát Watame: ``Ừ, cũng chỉ là một món ăn ngon thôi mà. Sao phải lo lắng thế? Cứ tận hưởng đi, Watame.''

Tôi cười trừ, vừa nhai miếng thịt vừa nói: ``Miễn là nó giúp em phục hồi tinh thần sau cú sốc `con cừu thách đấu' thì\dots\ cứ tận hưởng đi. Hôm nay em đã trải qua đủ chuyện rồi.''

Watame bỗng chốc khựng lại, như thể vừa nhận ra điều gì đó. Cô nhìn miếng thịt trong tay, rồi nhìn tôi, rồi nhìn Coco và Kanata. Nhưng rồi mùi vị thơm ngon của miếng thịt mà cô đang ăn đã nhanh chóng lấn át mọi suy nghĩ trong đầu cô, và cô tiếp tục ăn với vẻ sảng khoái.

Tôi đưa mắt nhìn sang ba người bạn của Watame đang cười đùa vui vẻ. Coco đang kể chuyện gì đó về một lần cô và Kanata phá hỏng bếp, còn Kanata thì đỏ mặt cãi lại. Watame vừa ăn vừa cười, mắt long lanh.

Nhưng rồi, tôi ngước mắt lên trần nhà, khẽ thở dài. Nghĩ đến việc mình vẫn còn hai thành viên nữa từ \textit{holoForce} chưa lên văn phòng, tôi bất giác thốt lên, giọng đầy chán nản: ``Haa\dots\ Xem ra sắp tới lại có chuyện rồi. Mới có ba thành viên mà đã đủ thứ: đột kích, bom, cừu thách đấu\dots\ Còn hai người nữa, không biết sẽ ra sao.''

Coco nghe tôi nói, bật cười: ``Kuon-paisen đừng lo, có lẽ hai người cuối cùng sẽ ngoan hơn bọn em nhiều.''

Kanata nhìn cô, giọng đầy nghi ngờ: ``Bà nghĩ thế á? Có khi họ còn tệ hơn đấy.''

Watame đặt miếng thịt xuống, nhìn tôi với ánh mắt đầy tò mò: ``Anh có linh cảm gì về hai người tiếp theo không?''

Tôi chỉ nhún vai, giọng đầy bí ẩn: ``Anh không biết, nhưng anh có cảm giác cái tên `na-nora\~{}' sẽ còn ám ảnh anh một thời gian dài.''

Ở đâu đó, một tiếng ``na-nora\~{}'' đầy nũng nịu vang lên khe khẽ như một điềm báo không mấy sáng sủa sắp sửa xảy ra với tôi.

\begin{center}
    {\Large \abv{Bonus}}
\end{center}

\begin{figure}[H]
    \centering
    \includegraphics[width=12.5cm]{../../../../Image/Book5/5-7m.png}
\end{figure}

\end{document}